\documentclass[preprint,12pt]{elsarticle}

\usepackage[T1]{fontenc}
\usepackage{amsmath,amssymb,amsthm,mathtools}
\usepackage{booktabs,longtable,array}
\usepackage{microtype}
\usepackage[hidelinks]{hyperref}
\usepackage{enumitem}
\usepackage{xcolor}

\newtheorem{theorem}{Theorem}[section]
\newtheorem{lemma}[theorem]{Lemma}
\newtheorem{proposition}[theorem]{Proposition}
\newtheorem{corollary}[theorem]{Corollary}
\newtheorem{certificate}[theorem]{Lemma}
\theoremstyle{definition}
\newtheorem{definition}[theorem]{Definition}
\newtheorem{remark}[theorem]{Remark}

\newcommand{\Z}{\mathbb Z}
\newcommand{\Q}{\mathbb Q}
\newcommand{\R}{\mathbb R}
\newcommand{\N}{\mathbb N}
\newcommand{\dd}{\,\mathrm d}
\newcommand{\e}{\mathrm e}
\DeclareMathOperator{\lcm}{lcm}
\newcommand{\eps}{\varepsilon}
\newcommand{\abs}[1]{\left|#1\right|}

\journal{Journal of Number Theory}

\hypersetup{
  pdftitle={Improved irrationality-exponent bounds for log 3 and log base 2 of 3},
  pdfauthor={Rohit Kumar Jha},
  pdfsubject={Diophantine approximation and irrationality exponents},
  pdfkeywords={irrationality exponent, logarithms, computer-assisted proof,
    exact real-root isolation}
}

\begin{document}
\begin{frontmatter}

\title{Improved irrationality-exponent bounds for
$\boldsymbol{\log 3}$ and $\boldsymbol{\log_2 3}$}

\author{Rohit Kumar Jha}

\begin{abstract}
We prove that, for all sufficiently large
$H=\max(|h_1|,|h_2|)$ and all integers $h_0,h_1,h_2$ with
$(h_1,h_2)\ne(0,0)$,
\[
 |h_0+h_1\log2+h_2\log3|>H^{-4.104104}.
\]
Consequently,
$\mu(\log3)\le5.104104$ and
$\mu(\log3/\log2)\le5.104104$, improving the previously published
bound $5.116201$.  The proof constructs two positive integer linear forms
with a common coefficient in $\log(4/3)$ and $\log(9/8)$ from a product of
twenty-one symmetrized integer polynomials with exact rational weights.
Coefficient-content identities provide integrality, the Laplace principle
gives two distinct decay rates, and Cauchy's estimate controls the common
coefficient.  The finite inequalities are certified in exact rational
arithmetic by real-root isolation, Bernstein-basis sign certificates, and
rational logarithm enclosures.  A limit-superior two-form criterion then
yields the homogeneous estimate and both irrationality-exponent bounds.
\end{abstract}

\begin{keyword}
irrationality exponent \sep logarithms \sep integer linear forms \sep
symmetrized polynomials \sep computer-assisted proof \sep exact real-root isolation
\MSC[2020] 11J82 \sep 11Y60
\end{keyword}

\end{frontmatter}

\section{Introduction and statement of the result}

Throughout, $\log$ denotes the natural logarithm.  For an irrational real
number $\xi$, its irrationality exponent is
\[
\begin{aligned}
 \mu(\xi)=\inf\bigl\{\mu>0:\;&
 0<|\xi-p/q|<q^{-\mu}\ \text{has only finitely many solutions}\\
 & (p,q)\in\mathbb Z\times\mathbb Z_{\ge1}\bigr\}.
\end{aligned}
\]
Dirichlet's theorem gives $\mu(\xi)\ge2$ for every irrational $\xi$.
Obtaining an explicit upper bound for a particular transcendental constant is
a different problem: one must construct rational approximations and control,
simultaneously and effectively, both their denominators and their errors.

The two numbers considered here are connected by the linear form
\[
 q\log3-p\log2=\log\frac{3^q}{2^p}.
\]
Thus rational approximation to $\log3/\log2=\log_2 3$ measures how closely
powers of $2$ and $3$ can approach one another on a logarithmic scale.  A
lower bound for arbitrary integer forms in $1,\log2,\log3$ is stronger than
separate irrationality-exponent estimates and yields both of them by simple
specialization.  Such bounds also have applications outside number theory:
Draper and Saad use the exponent $5.116201$ in a space--entropy lower bound
for exact sampling from the Bernoulli distribution with parameter $1/3$
\cite{DraperSaad2026}.

Effective bounds for logarithms of rational numbers are commonly obtained
from definite integrals whose partial-fraction expansions are rational linear
forms in logarithms; see Zudilin's survey \cite{Zudilin2004}.  The quality of
such a construction is governed by three competing quantities: the
exponential decay of the integral, the growth of the logarithm coefficient,
and the arithmetic cost of clearing denominators.  Rhin's simultaneous
construction gave $\mu(\gamma)<8.616$ for every nonzero
$\gamma\in\mathbb Q\log2+\mathbb Q\log3$ \cite{Rhin1987}.  For $\log3$,
Salikhov
obtained
$\mu(\log3)\le5.125\ldots$ \cite{Salikhov2007}, and Wu and Wang reduced the
bound to $5.1163051$ using an arithmetic optimization with symmetrized linear
factors \cite{WuWang2014}.  Bondareva, Luchin, and Salikhov added a quadratic
symmetrized factor and proved the homogeneous exponent $4.116201$
\cite{BLS2018}.  In particular, their result gives
\begin{equation}\label{eq:previous-bound}
 \mu(\log3)\le5.116201,
 \qquad
 \mu(\log3/\log2)\le5.116201.
\end{equation}

The present construction enlarges the symmetrized factor family to twenty-one
polynomials of degrees between one and five and uses an exactly defined
positive rational weight vector.  The local coefficient-content balances are
retained, while the resulting
coefficient-growth-to-decay quotient is reduced from $4.116201$ to below
$4.104104$.  This gives the following strict
improvement of \eqref{eq:previous-bound}.

\begin{theorem}\label{thm:main}
Let $h_0,h_1,h_2\in\Z$, with $(h_1,h_2)\ne(0,0)$, and set
$H=\max(|h_1|,|h_2|)$.  For all sufficiently large $H$,
\[
 \abs{h_0+h_1\log2+h_2\log3}>H^{-4.104104}.
\]
\end{theorem}

\begin{corollary}\label{cor:irrationality}
One has
\[
 \mu(\log3)\le5.104104,
 \qquad
 \mu\!\left(\frac{\log3}{\log2}\right)\le5.104104.
\]
\end{corollary}

Here is the structure of the proof.  With
\[
 \theta_1=\log(4/3),\qquad \theta_2=\log(9/8),
\]
the symmetrized rational function produces two positive integrals over
$[35,40]$ and $[40,42]$.  After an exact coefficient-content calculation and
partial fractions, multiplication by a common clearing factor turns them into
integer linear forms
\[
 q_n\theta_1-p_{1,n},\qquad q_n\theta_2-p_{2,n}
\]
with the same coefficient $q_n$.  Positivity and the Laplace principle reduce
their exponential decay to maxima of a single real function on two adjacent
intervals.  Cauchy's coefficient estimate reduces the growth of $q_n$ to a
maximum on a complex circle.  Exact root isolation, Bernstein-basis signs,
and rational logarithm enclosures then prove that the two decay rates are
distinct and that the quotient of coefficient growth by the smaller decay
rate is less than
\[
 4.104103101316511697013643<4.104104.
\]
Finally, the unimodular relations
\[
 \log(4/3)=2\log2-\log3,
 \qquad
 \log(9/8)=-3\log2+2\log3
\]
transfer the resulting homogeneous estimate to $\log2$ and $\log3$.

Two features require explicit treatment.  First, the Cauchy argument gives a
limit superior for the common coefficient rather than a full exponential
limit.  Section~2 therefore proves the precise limit-superior version of
Hata's two-form criterion \cite[Lemma~2.1]{Hata1993} needed here.  Second,
the optimization inequalities are too large for useful printed expansion but
form a finite certificate over $\mathbb Q$.  The main proof states exactly the
certified polynomial and interval assertions it uses; their exact reproduction
is described in~\ref{app:certificate}, and the complete certificate accompanies
the paper as supplementary material.  No floating-point comparison is used in
a theorem-facing decision.

Section~2 establishes the Diophantine criterion.  Sections~3 and~4 define the
construction and prove integrality.  Sections~5--7 determine the two decay
rates and bound the common coefficient, and Section~8 completes the proof.
The factor data are listed in~\ref{app:factors}, while the finite certificate
and its reproduction are recorded in~\ref{app:certificate}.

\section{A two-linear-form criterion}

The next lemma is a limit-superior variant of Hata's two-form criterion
\cite[Lemma~2.1]{Hata1993}.  We include the proof because replacing a limit by
a limit superior is not merely notational: possible cancellation between the
two small forms must be controlled uniformly in their integer coefficients.

\begin{lemma}[two-form criterion]\label{lem:two-form}
Let $\theta_1,\theta_2\in\R$.  Suppose that integers
$q_n,p_{1,n},p_{2,n}$ satisfy
\[
 \eps_{j,n}=q_n\theta_j-p_{j,n}\ne0\qquad(j=1,2)
\]
for all sufficiently large $n$, and suppose that
\begin{equation}\label{eq:two-form-hypotheses}
 \lim_{n\to\infty}\frac1n\log\abs{\eps_{j,n}}=-\tau_j
 \quad(j=1,2),
 \qquad
 \limsup_{n\to\infty}\frac1n\log\abs{q_n}\le\lambda,
\end{equation}
where $\tau_1,\tau_2>0$, $\tau_1\ne\tau_2$, and $\lambda\ge0$.
Put $\tau=\min(\tau_1,\tau_2)$.  For every $\nu>\lambda/\tau$ there is
$H_0(\nu)$ such that
\[
 \abs{h_0+h_1\theta_1+h_2\theta_2}>H^{-\nu}
\]
whenever $h_0,h_1,h_2\in\Z$,
$H=\max(|h_1|,|h_2|)\ge H_0(\nu)$, and
$(h_1,h_2)\ne(0,0)$.
\end{lemma}

\begin{proof}
Interchanging the two coordinates if necessary, assume
$0<\tau_1<\tau_2$.  Write $d=\tau_2-\tau_1$ and $T=\tau_2$.
Fix $\nu>\lambda/\tau_1$.  Choose $c>0$ and then $\eta>0$ so small that
\begin{gather}
 dc>2\eta(2+3c),
 \qquad
 c\tau_1>\eta(2+2c),                                      \label{eq:eta-conditions}\\
 \frac{\Gamma}{\tau_1-\eta}<\nu,                           \label{eq:gamma-choice}\\
 \Gamma=(\lambda+\eta)(1+2c)+2cT+(2+2c)\eta.              \label{eq:gamma}
\end{gather}
Such a choice is possible: take $c\downarrow0$ and then
$\eta=o(c)$; the quotient in \eqref{eq:gamma-choice} tends to
$\lambda/\tau_1$.

After increasing $N$, the hypotheses give, for $n\ge N$ and $j=1,2$,
\begin{equation}\label{eq:two-form-exp-bounds}
 \e^{-(\tau_j+\eta)n}
 \le\abs{\eps_{j,n}}
 \le\e^{-(\tau_j-\eta)n},
 \qquad
 0<\abs{q_n}\le\e^{(\lambda+\eta)n}.
\end{equation}
The eventual nonvanishing of $q_n$ follows as well: if $q_n=0$ when both
errors have absolute value less than one, then
$\eps_{j,n}=-p_{j,n}$ cannot be a nonzero integer.

Fix $(h_0,h_1,h_2)$ and put
\[
 S_m=|h_1|\abs{\eps_{1,m}}+|h_2|\abs{\eps_{2,m}},
 \qquad
 E_m=h_1\eps_{1,m}+h_2\eps_{2,m}.
\]
Let $n\ge N+1$ be the least integer for which $S_n<1/2$.
This integer exists.  Since one of $|h_1|$ and $|h_2|$ equals $H$,
\[
 S_m\ge H\min\bigl(|\eps_{1,m}|,|\eps_{2,m}|\bigr).
\]
For every fixed $M$, the nonzero errors at the finitely many indices
$N+1,\ldots,M$ therefore show that $S_m\ge1/2$ throughout that range once
$H$ is sufficiently large.  Thus $n\to\infty$ uniformly with $H$; in
particular, $n>N+1$ and $S_{n-1}\ge1/2$ for sufficiently large $H$.
By \eqref{eq:two-form-exp-bounds},
\[
 \frac12\le S_{n-1}
 \le2H\e^{-(\tau_1-\eta)(n-1)},
\]
so
\begin{equation}\label{eq:n-height}
 n\le1+\frac{\log(4H)}{\tau_1-\eta}.
\end{equation}

Set
\[
 k_1=\lceil(1+c)n\rceil,
 \qquad
 k_2=\lceil(1+2c)n\rceil.
\]
The inequality $S_n<1/2$ and the lower bounds in
\eqref{eq:two-form-exp-bounds} imply
\[
 |h_j|<\frac12\e^{(\tau_j+\eta)n}\qquad(j=1,2).
\]
Consequently, for $k\in\{k_1,k_2\}$,
\[
 |h_j\eps_{j,k}|
 \le\frac12
 \exp\!\bigl(-\tau_j(k-n)+\eta(n+k)\bigr).
\]
The second inequality in \eqref{eq:eta-conditions} now shows that
$S_{k_1},S_{k_2}<1/2$ once $H$, and hence $n$, is sufficiently large.

Suppose first that $h_1h_2\ne0$, and define
\[
 R_k=\frac{|h_1\eps_{1,k}|}{|h_2\eps_{2,k}|}.
\]
The coefficient quotient cancels from $R_{k_2}/R_{k_1}$.  From
\eqref{eq:two-form-exp-bounds},
\[
 \frac{R_{k_2}}{R_{k_1}}
 \ge
 \exp\!\bigl((d-2\eta)k_2-(d+2\eta)k_1\bigr)>4
\]
for all sufficiently large $n$, by the first inequality in
\eqref{eq:eta-conditions}.  Hence at least one
$k\in\{k_1,k_2\}$ satisfies
$R_k\notin[1/2,2]$.  Choose such a $k$.  If one of $h_1,h_2$ is zero,
choose $k=k_1$ instead.

Since $S_{n-1}\ge1/2$, there is an index $j_0\in\{1,2\}$ with
\[
 |h_{j_0}\eps_{j_0,n-1}|\ge\frac14.
\]
Using first the upper and then the lower error bounds in
\eqref{eq:two-form-exp-bounds}, and using
$k\le(1+2c)n+1$, gives
\begin{align*}
 |h_{j_0}\eps_{j_0,k}|
 &\ge\frac14
 \exp\!\bigl((\tau_{j_0}-\eta)(n-1)
             -(\tau_{j_0}+\eta)k\bigr)\\
 &\ge\frac{\e^{-2T}}4
 \exp\!\bigl(-(2cT+(2+2c)\eta)n\bigr).
\end{align*}
If both coefficients are nonzero, the choice
$R_k\notin[1/2,2]$ implies that one of the two summands in $E_k$ has
more than twice the absolute value of the other.  If one coefficient is zero,
there is only one summand.  In either case,
\begin{equation}\label{eq:E-lower}
 |E_k|\ge\frac{\e^{-2T}}8
 \exp\!\bigl(-(2cT+(2+2c)\eta)n\bigr).
\end{equation}

Put
\[
 \Lambda=h_0+h_1\theta_1+h_2\theta_2,
 \qquad
 A_k=q_k h_0+p_{1,k}h_1+p_{2,k}h_2\in\Z.
\]
Then $q_k\Lambda=A_k+E_k$ and $|E_k|\le S_k<1/2$.
If $A_k\ne0$, then
$|\Lambda|>1/(2|q_k|)$.  If $A_k=0$, then
$|\Lambda|=|E_k|/|q_k|$, and \eqref{eq:E-lower} applies.
Together with \eqref{eq:two-form-exp-bounds}, both cases give a constant
$C>0$, independent of the coefficients, such that
\[
 |\Lambda|\ge C\e^{-\Gamma n}.
\]
Using \eqref{eq:n-height}, this is at least a positive constant times
$H^{-\Gamma/(\tau_1-\eta)}$.  By \eqref{eq:gamma-choice}, it exceeds
$H^{-\nu}$ for all sufficiently large $H$.  This proves the lemma.
\end{proof}

\section{The symmetrized integral construction}
\label{sec:construction}

The numerator will be invariant under $x\mapsto70-x$.  It is therefore
convenient to put
\[
 t=(x-35)^2.
\]
The twenty-one integer polynomials $P_k(t)$ are listed, together with their
local coefficient profiles, in Table~\ref{tab:factors} of~\ref{app:factors}.
Their degrees are between one and five.

For $p\in\{2,3,5,7\}$, define
\begin{equation}\label{eq:profile-definition}
 \tau_p(P_k)=
 \min_j v_p\!\left([y^j]P_k((840y-35)^2)\right),
\end{equation}
where zero coefficients are omitted from the minimum.  Thus
$p^{\tau_p(P_k)}$ is the exact $p$-part of the coefficient content relevant
to the scaling used below.

\subsection{Exact weights}

Set
\begin{equation}\label{eq:z0}
 z_0=26\frac{-899+60i}{901},
 \qquad
 t_0=(z_0-35)^2
 =\frac{3012564681-171316080i}{811801}.
\end{equation}
In particular, $|z_0|=26$.  Direct exact substitution gives
$P_k(t_0)\ne0$ for every $k$.  For $1\le k\le21$, put
\begin{equation}\label{eq:ellk}
 \ell_k=
 2(z_0-35)\frac{P_k'(t_0)}{P_k(t_0)}\in\Q(i).
\end{equation}

The weights $\alpha_k\in\Q_{>0}$ are defined exactly as follows.  For
$7\le k\le21$, set $\alpha_k=u_k/10^{12}$, where
\begin{center}
\small
\begin{tabular}{crrrr}
\toprule
$k$&7&8&9&10\\
\midrule
$u_k$&119873656&878336597&101691070&103130300\\
\bottomrule
\end{tabular}
\qquad
\begin{tabular}{crrrr}
\toprule
$k$&11&12&13&14\\
\midrule
$u_k$&219953072&757245952&781996563&232727309\\
\bottomrule
\end{tabular}

\medskip
\begin{tabular}{crrrr}
\toprule
$k$&15&16&17&18\\
\midrule
$u_k$&630848726&330420926&232427174&314945936\\
\bottomrule
\end{tabular}
\qquad
\begin{tabular}{crrr}
\toprule
$k$&19&20&21\\
\midrule
$u_k$&103903210&14385085&9664407\\
\bottomrule
\end{tabular}
\end{center}
The first six weights are the unique rational solution of
\begin{align}
 \sum_{k=1}^{21}\alpha_k\deg P_k&=\frac32,
 \label{eq:degree-balance}\\
 \sum_{k=1}^{21}\alpha_k\tau_7(P_k)&=2,
 &\sum_{k=1}^{21}\alpha_k\tau_5(P_k)&=2,
 \label{eq:profile75}\\
 \sum_{k=1}^{21}\alpha_k\tau_3(P_k)&=1,
 \label{eq:profile3}\\
 \sum_{k=1}^{21}\alpha_k\ell_k
 &=\frac1{z_0}-\frac1{70-z_0}.
 \label{eq:complex-weight-equation}
\end{align}
The last line means equality of real and imaginary parts, so the system has
six rational equations.  Exact Gaussian elimination shows that the resulting
$6\times6$ coefficient matrix is nonsingular and that all six weights are
positive.
For orientation,
\begin{center}
\small
\begin{tabular}{r@{\quad}l r@{\quad}l r@{\quad}l}
\toprule
$k$&$\alpha_k$&$k$&$\alpha_k$&$k$&$\alpha_k$\\
\midrule
1 & 0.492767795602391 & 8 & 0.000878336597000 & 15 & 0.000630848726000 \\
2 & 0.497597698986041 & 9 & 0.000101691070000 & 16 & 0.000330420926000 \\
3 & 0.484860590979781 & 10 & 0.000103130300000 & 17 & 0.000232427174000 \\
4 & 0.000043741872954 & 11 & 0.000219953072000 & 18 & 0.000314945936000 \\
5 & 0.000129471059917 & 12 & 0.000757245952000 & 19 & 0.000103903210000 \\
6 & 0.000176164327825 & 13 & 0.000781996563000 & 20 & 0.000014385085000 \\
7 & 0.000119873656000 & 14 & 0.000232727309000 & 21 & 0.000009664407000 \\
\bottomrule
\end{tabular}
\end{center}

Let $D$ be the least common multiple of the denominators of the
$\alpha_k$, and write
\[
 W_k=D\alpha_k\in\Z_{>0}.
\]
The integer $D$ has $1061$ decimal digits.  Its exact value and all $W_k$
are part of the supplementary certificate.  We let $n$ tend to infinity
through the positive multiples of $2D$.  Then every $\alpha_k n$ is an even
integer.

The numerical search that led to these parameters is not part of the proof.
Its provenance and the subsequent exact rationalization are recorded
in~\ref{app:certificate} and in the supplementary material.  No
optimality is claimed for the factor family.

\section{Arithmetic and integer linear forms}

Put
\begin{equation}\label{eq:An-Rn}
 A_n(x)=\prod_{k=1}^{21}
 P_k((x-35)^2)^{\alpha_k n},
 \qquad
 R_n(x)=
 \frac{A_n(x)}{x^{n+1}(70-x)^{n+1}}.
\end{equation}
The degree balance gives $\deg_x A_n=3n$.  The numerator has integer
coefficients, and all of its factor exponents are even.  Consequently
$R_n$ is nonnegative on the two real intervals used below.

\subsection{The scaled numerator}

Set
\begin{equation}\label{eq:a-clearing}
 s_2=\sum_{k=1}^{21}\alpha_k\tau_2(P_k),
 \qquad
 a=4-s_2.
\end{equation}
The exact weights give
\begin{equation}\label{eq:a-value}
 a=0.43339168265643340846553249867665\ldots>0.
\end{equation}
Because $n$ is a multiple of $2D$, $an$ is a nonnegative integer.

\begin{lemma}[integer scaled numerator]\label{lem:scaled-integer}
For every admissible $n$,
\begin{equation}\label{eq:Bn}
 B_n(y)=\frac{2^{an}A_n(840y)}{58800^n}
\end{equation}
belongs to $\Z[y]$.
\end{lemma}

\begin{proof}
For each $k$, definition \eqref{eq:profile-definition} says that every
coefficient of $P_k((840y-35)^2)$ is divisible by
\[
 2^{\tau_2(P_k)}3^{\tau_3(P_k)}
 5^{\tau_5(P_k)}7^{\tau_7(P_k)}.
\]
After raising the factors to the integer powers $\alpha_k n$ and multiplying,
\eqref{eq:profile75}--\eqref{eq:profile3} supply respectively
$7^{2n}$, $5^{2n}$, and $3^n$.  Definition \eqref{eq:a-clearing} supplies
$2^{4n}$ after multiplication by $2^{an}$.  Since
\[
 58800=2^4\cdot3\cdot5^2\cdot7^2,
\]
the quotient in \eqref{eq:Bn} has integer coefficients.
\end{proof}

\subsection{Partial fractions and denominator clearing}

Symmetry under $x\mapsto70-x$ gives
\begin{equation}\label{eq:partial-fraction}
 R_n(x)=Q_{n-2}(x)+
 \sum_{i=1}^{n+1}a_{i,n}
 \left(\frac1{x^i}+\frac1{(70-x)^i}\right),
\end{equation}
where $Q_{n-2}\in\Z[x]$.  Indeed, the polynomial denominator in
\eqref{eq:An-Rn} has leading coefficient $\pm1$, so polynomial division
preserves integrality, and symmetry identifies the two principal parts.

\begin{lemma}[partial-fraction denominators]\label{lem:ai}
For $1\le i\le n+1$ there is an integer $M_{i,n}$ such that
\begin{equation}\label{eq:ai}
 a_{i,n}=\frac{2^{-an}840^{i-1}}{70}M_{i,n}.
\end{equation}
\end{lemma}

\begin{proof}
At $x=0$,
\[
 a_{i,n}=[x^{n+1-i}]
 \frac{A_n(x)}{(70-x)^{n+1}}.
\]
Put $j=n+1-i$ and substitute $x=840y$.  By
Lemma~\ref{lem:scaled-integer},
\[
\begin{aligned}
 a_{i,n}
 &=840^{-j}\frac{2^{-an}58800^n}{70^{n+1}}
 [y^j]B_n(y)(1-12y)^{-n-1}\\
 &=\frac{2^{-an}840^{n-j}}{70}
 [y^j]B_n(y)(1-12y)^{-n-1}.
\end{aligned}
\]
Now $n-j=i-1$ and $(1-12y)^{-n-1}\in\Z[[y]]$.
\end{proof}

Let
\begin{equation}\label{eq:Cn}
 \Delta_n=\lcm(1,2,\ldots,n),
 \qquad
 C_n=70\,2^{an}\Delta_n,
\end{equation}
and define
\begin{equation}\label{eq:integrals}
 I_{1,n}=\int_{35}^{40}R_n(x)\dd x,
 \qquad
 I_{2,n}=\int_{40}^{42}R_n(x)\dd x.
\end{equation}

\begin{proposition}[two positive integer forms]\label{prop:forms}
There are integers $q_n,p_{1,n},p_{2,n}$ such that
\begin{align}
 C_n I_{1,n}&=q_n\log\frac43-p_{1,n},\label{eq:form1}\\
 C_n I_{2,n}&=q_n\log\frac98-p_{2,n},\label{eq:form2}
\end{align}
where the common coefficient is $q_n=C_n a_{1,n}\in\Z$.  Both left-hand
sides are strictly positive.
\end{proposition}

\begin{proof}
For integers $0<u<v<70$,
\begin{align}
\int_u^v\left(\frac1x+\frac1{70-x}\right)\dd x
 &=\log\frac{v(70-u)}{u(70-v)},\label{eq:elementary-log}\\
\int_u^v\left(\frac1{x^i}+\frac1{(70-x)^i}\right)\dd x
 &=\frac{u^{1-i}-v^{1-i}+(70-v)^{1-i}-(70-u)^{1-i}}{i-1}
 \label{eq:elementary-rational}
\end{align}
for $i\ge2$.  The logarithm in \eqref{eq:elementary-log} is
$\log(4/3)$ for $(u,v)=(35,40)$ and $\log(9/8)$ for
$(u,v)=(40,42)$.

Every endpoint in \eqref{eq:elementary-rational} belongs to
$\{28,30,35,40,42\}$ and divides $840$.  By \eqref{eq:ai},
multiplication by $70\,2^{an}$ clears its $(i-1)$st power, while
$\Delta_n$ clears the factor $i-1$.  The polynomial part has integer
coefficients and degree $n-2$, so $\Delta_n$ also clears all denominators in
its integral between integer endpoints.  This proves
\eqref{eq:form1}--\eqref{eq:form2}.  Positivity follows because $R_n$ is
nonnegative and not identically zero on either interval.
\end{proof}

Since $\log\Delta_n=\psi(n)$, the prime number theorem in the form
$\psi(n)\sim n$ gives
\[
 \lim_{n\to\infty}\frac1n\log\Delta_n=1
\]
\cite[Ch.~6]{MontgomeryVaughan2007}.  Hence, along the admissible values of
$n$,
\begin{equation}\label{eq:cdef}
 \lim_{n\to\infty}\frac1n\log C_n
 =c:=1+a\log2.
\end{equation}

\section{Exact logarithm enclosures}\label{sec:log-enclosure}

\begin{lemma}\label{lem:log-enclosure}
For $0\le z<1$ and $N\ge1$,
\[
 2\sum_{j=0}^{N-1}\frac{z^{2j+1}}{2j+1}
 \le\log\frac{1+z}{1-z}
 \le
 2\sum_{j=0}^{N-1}\frac{z^{2j+1}}{2j+1}
 +\frac{2z^{2N+1}}{(2N+1)(1-z^2)}.
\]
Consequently every logarithm of a positive rational number can be enclosed
between explicit rational numbers.
\end{lemma}

\begin{proof}
Use the positive series
$\operatorname{atanh}z=\sum_{j\ge0}z^{2j+1}/(2j+1)$ and bound every omitted
denominator below by $2N+1$.  For a positive rational $y$, choose the exact
integer $k$ with $2^k\le y<2^{k+1}$ and write
\[
 \log y=k\log2+
 \log\frac{1+u}{1-u},
 \qquad
 u=\frac{y/2^k-1}{y/2^k+1},\qquad 0\le u<\frac13.
\]
The same series with $z=1/3$ encloses
$\log2=\log((1+z)/(1-z))$, so every term in the displayed decomposition has
explicit rational lower and upper bounds.
\end{proof}

\section{The two real decay rates}

For $0\le t\le49$, define
\begin{equation}\label{eq:g}
 g(t)=\frac{\prod_{k=1}^{21}\abs{P_k(t)}^{\alpha_k}}{1225-t}.
\end{equation}
This is a continuous nonnegative function.  On the two integration
intervals,
\begin{equation}\label{eq:Rn-g}
 R_n(x)=
 \frac{g((x-35)^2)^n}{1225-(x-35)^2}.
\end{equation}
Put
\begin{equation}\label{eq:M12}
 M_1=\log\max_{0\le t\le25}g(t),
 \qquad
 M_2=\log\max_{25\le t\le49}g(t).
\end{equation}

\begin{lemma}[elementary Laplace principle]\label{lem:laplace}
Let $f$ be continuous and nonnegative, but not identically zero, on a
nondegenerate compact interval $K$, and let $w$ be continuous and strictly
positive there.  Then
\[
 \lim_{n\to\infty}\frac1n
 \log\int_K f(x)^n w(x)\dd x
 =\log\max_K f.
\]
\end{lemma}

\begin{proof}
The upper bound follows from
$\int_K f^n w\le(\max_K f)^n\int_K w$.  For the lower bound, choose a point at
which $f$ attains its positive maximum.  For every $\delta>0$, the function
$f$ is at least $\max_K f-\delta$ on a one-sided or two-sided neighborhood
of that point, and $w$ has a positive lower bound there.  Let $n\to\infty$
and then $\delta\downarrow0$.
\end{proof}

Away from the zeros of the $P_k$,
\[
 \frac{\dd}{\dd t}\log g(t)
 =\sum_{k=1}^{21}\alpha_k\frac{P_k'(t)}{P_k(t)}
 +\frac1{1225-t}.
\]
Let $P(t)=\prod_{k=1}^{21}P_k(t)$, and let $H(t)$ be the primitive part,
with positive content removed, of
\begin{equation}\label{eq:H}
 (1225-t)\sum_{k=1}^{21}W_k P_k'(t)\frac{P(t)}{P_k(t)}
 +DP(t).
\end{equation}
It is an integer polynomial of degree $85$.

\begin{certificate}[real critical points]\label{cert:real}
The polynomial $H$ is squarefree and coprime to $P$.  It has exactly $63$
roots in $[0,49]$: $41$ in $(0,25)$ and $22$ in $(25,49)$.  Exact rational
root isolation and outward rational logarithm bounds give
\begin{align}
 -2.722957495668415052959634
 &<M_1<
 -2.722957495668409711051756,
 \label{eq:M1-bounds}\\
 -2.722977562389619550059089
 &<M_2<
 -2.722977562389619502813003.
 \label{eq:M2-bounds}
\end{align}
In particular,
\begin{equation}\label{eq:M-separation}
 M_1-M_2>0.000020066721204449853370212.
\end{equation}
\end{certificate}

\begin{proof}
This is a finite calculation over $\Q$, recorded in~\ref{app:certificate}.
After clearing content, exact Euclidean
algorithms prove that $H$ is squarefree and coprime to $P$, and the
Vincent--Akritas--Strzebo\'nski continued-fraction method isolates every real
root.
The endpoints $0,25,49$ and every other zero of a factor have $g(t)=0$;
therefore every positive maximum occurs at one of the isolated roots of
$H$.  Each factor is bounded on each rational isolating interval by exact
Horner interval arithmetic.  The logarithms are enclosed using
Lemma~\ref{lem:log-enclosure}.  Comparing the resulting rational intervals
gives \eqref{eq:M1-bounds}--\eqref{eq:M-separation}.  All $63$ isolating
intervals and their exact rational value bounds are included in the
supplementary certificate.
\end{proof}

By Lemma~\ref{lem:laplace}, Proposition~\ref{prop:forms}, and
\eqref{eq:cdef},
\begin{equation}\label{eq:error-rates}
 \lim_{\substack{n\to\infty\\2D\mid n}}
 \frac1n\log(C_n I_{j,n})
 =c+M_j=-\tau_j
 \qquad(j=1,2).
\end{equation}
Lemma~\ref{cert:real} shows that $\tau_2>\tau_1$, and the numerical
bounds in Section~\ref{sec:completion} show that both are positive.  Since
the two errors are positive, they are nonzero for every admissible $n$.

\section{Growth bound for the common coefficient}

Put
\[
 Q_k(z)=P_k((z-35)^2).
\]
The coefficient of the logarithms in Proposition~\ref{prop:forms} is
controlled by
\begin{equation}\label{eq:a1-coefficient}
 a_{1,n}=[z^n]\frac{A_n(z)}{(70-z)^{n+1}}.
\end{equation}
We apply Cauchy's estimate on the circle $|z|=26$.

Write $z=26\e^{i\theta}$ and $u=\cos\theta$.  Define the integer
polynomials
\begin{equation}\label{eq:circle-factors}
 A_k^\circ(u)=|Q_k(26\e^{i\theta})|^2
\end{equation}
and
\begin{equation}\label{eq:B-circle}
 B^\circ(u)=|26\e^{i\theta}(70-26\e^{i\theta})|^2
 =676(5576-3640u).
\end{equation}
Set
\begin{equation}\label{eq:F-circle}
 F(u)=\frac12\sum_{k=1}^{21}\alpha_k\log A_k^\circ(u)
 -\frac12\log B^\circ(u),
 \qquad -1\le u\le1,
\end{equation}
and let $L=\max_{[-1,1]}F$.

Let $\mathcal A(u)=\prod_k A_k^\circ(u)$.  The primitive integer numerator
of $2D F'(u)$ is the primitive part, with positive content removed, of
\begin{equation}\label{eq:J}
 B^\circ(u)\sum_{k=1}^{21}
 W_k(A_k^\circ)'(u)\frac{\mathcal A(u)}{A_k^\circ(u)}
 -D(B^\circ)'(u)\mathcal A(u).
\end{equation}
Denote it by $J(u)$.

\begin{certificate}[circle maximum]\label{cert:circle}
Every $A_k^\circ$ is strictly positive on $[-1,1]$.  The polynomial $J$
has degree $170$ and factors as
\begin{equation}\label{eq:J-factor}
 J(u)=(901u+899)K(u),
\end{equation}
where $K$ has degree $169$ and is strictly negative on $[-1,1]$.
Consequently $F$ has its unique maximum at
\[
 u_0=-\frac{899}{901}.
\]
Exact outward logarithm bounds at $u_0$ give
\begin{equation}\label{eq:L-bounds}
 4.537901075598406403746477
 <L<
 4.537901075598406404986681.
\end{equation}
\end{certificate}

\begin{proof}
Map $u=2v-1$ from $[-1,1]$ to $[0,1]$.  If
$p(v)=\sum_{j=0}^n a_j v^j$, its degree-$n$ Bernstein coefficients are
\begin{equation}\label{eq:bernstein}
 b_k=\sum_{j=0}^k a_j
 \frac{\binom{k}{j}}{\binom{n}{j}}
 \qquad(0\le k\le n).
\end{equation}
Exact conversion by \eqref{eq:bernstein} gives positive Bernstein
coefficients for every $A_k^\circ(2v-1)$ and negative Bernstein coefficients
for $K(2v-1)$.  Hence every circle factor is positive on $[-1,1]$, and the
sign of $J$ is positive to the left of $u_0$ and negative to its right.
The same calculation gives $J'(u_0)<0$.  The complete coefficient and sign
certificate is described in~\ref{app:certificate}.  Finally,
Lemma~\ref{lem:log-enclosure} gives \eqref{eq:L-bounds}.
\end{proof}

\begin{proposition}[Cauchy growth bound]\label{prop:cauchy}
Along the admissible values of $n$,
\begin{equation}\label{eq:q-limsup}
 \limsup_{\substack{n\to\infty\\2D\mid n}}
 \frac1n\log|q_n|
 \le c+L.
\end{equation}
\end{proposition}

\begin{proof}
From \eqref{eq:a1-coefficient} and Cauchy's coefficient estimate,
\[
\begin{aligned}
 |a_{1,n}|
 &\le
 \max_{|z|=26}
 \frac{|A_n(z)|}{26^n|70-z|^{n+1}}\\
 &\le\frac1{44}
 \left(
 \max_{|z|=26}
 \frac{\prod_k|Q_k(z)|^{\alpha_k}}{|z(70-z)|}
 \right)^n
 =\frac1{44}\e^{nL}.
\end{aligned}
\]
Since $q_n=C_n a_{1,n}$, equation \eqref{eq:cdef} gives
\eqref{eq:q-limsup}.  If a coefficient vanishes, the corresponding
logarithm is interpreted as $-\infty$.  By \eqref{eq:error-rates}, both
positive errors lie in $(0,1)$ for every sufficiently large admissible $n$.
If $q_n=0$ at such an index, \eqref{eq:form1} would make the first error a
nonzero integer, a contradiction.  Thus $q_n$ is eventually nonzero.
\end{proof}

\section{Certified constants and completion}\label{sec:completion}

The exact rational logarithm enclosure gives
\begin{equation}\label{eq:c-bounds}
 1.300404222911437365598758
 <c<
 1.300404222911437366032151.
\end{equation}
Combining the outward upper bounds for $c$, $M_1$, and $L$ gives
\begin{align}
 \tau_1=-(c+M_1)
 &>1.422553272756972345019606,
 \label{eq:tau-final}\\
 \lambda:=c+L
 &<5.838305298509843771018831.
 \label{eq:lambda-final}
\end{align}
The rational endpoints in these enclosures satisfy, by exact integer
cross-multiplication,
\begin{equation}\label{eq:ratio-final}
 \frac{\lambda}{\tau_1}
 <4.104103101316511697013643
 <\frac{513013}{125000}
 =4.104104.
\end{equation}

For $m\ge1$, reindex the construction by $n_m=2Dm$.  Equations
\eqref{eq:error-rates} and \eqref{eq:q-limsup} then satisfy the hypotheses of
Lemma~\ref{lem:two-form}, with every exponential constant multiplied by
$2D$ and hence with the quotient in \eqref{eq:ratio-final} unchanged.
Apply the lemma to
\[
 \theta_1=\log\frac43=2\log2-\log3,
 \qquad
 \theta_2=\log\frac98=-3\log2+2\log3.
\]

The coefficient transformation is unimodular.  Explicitly, for arbitrary
$h_1,h_2\in\Z$, put
\begin{equation}\label{eq:unimodular}
 u=2h_1+3h_2,
 \qquad
 v=h_1+2h_2.
\end{equation}
Then
\[
 u\theta_1+v\theta_2=h_1\log2+h_2\log3.
\]
Moreover, if $H=\max(|h_1|,|h_2|)$ and
$H'=\max(|u|,|v|)$, then $H/5\le H'\le5H$.
Choose $\nu$ strictly between the first and second quantities in
\eqref{eq:ratio-final}.  Lemma~\ref{lem:two-form} gives
\[
 \abs{h_0+h_1\log2+h_2\log3}
 >(H')^{-\nu}\ge(5H)^{-\nu}.
\]
For all sufficiently large $H$, the last expression exceeds
$H^{-4.104104}$.  This proves Theorem~\ref{thm:main}.

\begin{proof}[Proof of Corollary~\ref{cor:irrationality}]
Taking $(h_1,h_2)=(0,q)$ and $h_0=-p$ in
Theorem~\ref{thm:main} gives
\[
 \abs{\log3-\frac pq}>q^{-5.104104}
\]
for all sufficiently large $q$ and every $p\in\Z$.

For $\xi=\log3/\log2$, take $(h_1,h_2)=(-p,q)$.  Whenever
$|\xi-p/q|<1$, one has
$\max(|p|,q)\le(1+|\xi|)q$.  The theorem therefore gives a positive
constant $C_\xi$ such that
\[
 \abs{\xi-\frac pq}
 =\frac{\abs{q\log3-p\log2}}{q\log2}
 >C_\xi q^{-5.104104}
\]
for all sufficiently large $q$.  For every $\delta>0$, the right-hand side
exceeds $q^{-(5.104104+\delta)}$ once $q$ is sufficiently large.  Thus the
defining approximation inequality has only finitely many solutions for every
exponent greater than $5.104104$.  Together with the first lower bound, this
proves both assertions.
\end{proof}

\appendix

\makeatletter
\@addtoreset{table}{section}
\makeatother
\renewcommand{\thetable}{\Alph{section}.\arabic{table}}

\section{Factors and valuation profiles}\label{app:factors}

In the third column, coefficients are ordered from the leading coefficient
to the constant term.  The historical ID records the column in the
exploratory factor pool.

\scriptsize
\begin{longtable}{r r >{\raggedright\arraybackslash\ttfamily}p{0.55\textwidth} rrrr}
\caption{The twenty-one symmetrized factors and their local coefficient
profiles.}\label{tab:factors}\\
\toprule
$k$&ID&\normalfont coefficients, high to low&$\tau_7$&$\tau_5$&$\tau_3$&$\tau_2$\\
\midrule
\endfirsthead
\multicolumn{7}{c}{\tablename\ \thetable\ (continued)}\\
\toprule
$k$&ID&\normalfont coefficients, high to low&$\tau_7$&$\tau_5$&$\tau_3$&$\tau_2$\\
\midrule
\endhead
1&0&[1, -49]&2&0&1&3\\
2&1&[1, -25]&0&2&1&4\\
3&2&[1, 0]&2&2&0&0\\
4&43&[38553, -19202708, 1088493350, -17941472500, 90075015625]&8&6&2&16\\
5&99&[463, -2770460, 172374650, -2731137500, 12867859375]&7&6&2&15\\
6&105&[1579, 4041716, -242861150, 3823592500, -18015003125]&8&5&1&16\\
7&124&[11, -38729, 834225, -4501875]&4&4&1&12\\
8&144&[16201, -64316693, 6078999290, -194217490250, 2344523978125, -9457876640625]&9&7&2&20\\
9&148&[3723, -14434700, 867361250, -13655687500, 64339296875]&7&7&0&16\\
10&159&[154073, -706525365, 69673778650, -2292537826250, 27833179828125, -110341894140625]&10&8&2&20\\
11&161&[13627, -789705, 11790625, -52521875]&5&5&1&12\\
12&163&[4259, -19924359, 1880315710, -60526208750, 754056559375, -3152625546875]&9&7&1&20\\
13&178&[8831, -33425731, 2301541270, -49661683750, 434933646875, -1351125234375]&8&7&2&20\\
14&189&[11881, -45475525, 3256944250, -74061846250, 685673078125, -2251875390625]&8&8&2&19\\
15&190&[31809, -138944253, 9336110826, -192597295450, 1574511273125, -4413675765625]&10&6&2&20\\
16&195&[85343, -320906467, 22689947350, -517088363750, 4799711546875, -15763127734375]&9&8&3&19\\
17&198&[21629, -86995629, 6252170386, -142667059850, 1329507230625, -4413675765625]&10&6&2&18\\
18&199&[359787, -1983981727, 196332531150, -6510548598750, 80617138984375, -331025682421875]&10&8&3&20\\
19&201&[14179, -60400263, 4236862910, -96560416750, 918765159375, -3152625546875]&9&7&3&20\\
20&202&[43951, -94114083, 9117300150, -301286483750, 3770282796875, -15763127734375]&9&8&2&19\\
21&208&[162521, -717601717, 70580396250, -2311361666250, 27833179828125, -110341894140625]&10&8&2&19\\
\bottomrule
\end{longtable}
\normalsize

\section{Computational certificate and reproducibility}
\label{app:certificate}

This appendix records the implementation of the finite calculations used in
Lemmas~\ref{cert:real} and~\ref{cert:circle} and in the
final comparison \eqref{eq:ratio-final}.  The analytic and arithmetic
reductions are proved in the main text; the role of the software is limited to
polynomial identities, real-root counts, rational interval inequalities, and
exact sign decisions.

\subsection{Origin and exact rationalization of the parameters}

The active factor support arose from a discretized optimization of a
semi-infinite linear program.  The surviving exploratory output had the
nonrigorous design value $5.1040765815\ldots$.  The submitted construction was
then rationalized by imposing the circle point \eqref{eq:z0}, prescribing a
separation margin between the two real maxima, fixing the last fifteen
weights with denominator $10^{12}$, and solving
\eqref{eq:degree-balance}--\eqref{eq:complex-weight-equation} exactly for the
first six weights.  The exploratory value is not used in the proof.  The
supplement contains the surviving design output and a numerical-to-exact
history so that this distinction can be audited.

\subsection{Primary exact certificate}

The program \texttt{verify\_log3\_bound\_5\_104104.py} uses Python integer
and rational arithmetic together with SymPy's exact polynomial algorithms
\cite{SymPy2017}.  Starting from the factor coefficients and the fifteen free
weights displayed in Section~\ref{sec:construction}, it
\begin{enumerate}[leftmargin=2em]
\item reconstructs the first six weights and checks their positivity and all
      six defining equations;
\item recomputes the twenty-one local coefficient profiles and all degree and
      content balances;
\item constructs the degree-$85$ polynomial $H$, proves squarefreeness and
      coprimality, and isolates its $63$ roots in $[0,49]$;
\item encloses every real critical value and proves the strict separation of
      $M_1$ and $M_2$;
\item constructs the degree-$170$ circle polynomial $J$, verifies
      $J=(901u+899)K$, and proves the required Bernstein signs;
\item encloses $L$, $c$, $\lambda$, and $\tau_1$ by rational intervals; and
\item proves \eqref{eq:ratio-final} by integer cross-multiplication.
\end{enumerate}
There is no floating-point branch or comparison, and the program refuses to
run when Python assertions are disabled.  Decimal output is diagnostic only.
The exact denominator, weight numerators, root intervals, value bounds, and
polynomial data are exported to
\texttt{certificate\_log3\_5\_104104.json}.

For identification of the reconstructed polynomials, the SHA-256 hashes of
their primitive high-to-low coefficient lists are
\begin{center}
\ttfamily\small
H: 91e83aa85d3e09743bd8664736bd4edf764a52b945bd46592b3bf98fc69749c3,\\
J: 71a2327cc73fa1909a5e90933b438db0c547d903cdc6e31021f2e82a0995ce92.
\end{center}

Every logarithm comparison implements Lemma~\ref{lem:log-enclosure}.  A
positive rational number is first reduced, by an exact power of two, to a
number $y$ satisfying $1\le y<2$;
the program then takes $24$ positive-series terms, appends the explicit
rational tail, and rounds outward to denominator $10^{18}$.

\subsection{Independent exact reproduction}

The second program, \texttt{independent\_exact\_audit.py}, uses only Python's
standard library and imports neither the primary verifier nor a
computer-algebra system.  It separately implements rational Gaussian
elimination, integer-polynomial arithmetic, logarithm enclosures, Bernstein
conversion, and the final comparison.  It reconstructs the coefficient lists
and hashes of $H$ and $J$.  Exact Descartes transformations prove that each of
the $63$ exported intervals contains one root of $H$ and that all $64$
complementary gaps in $(0,49)$ contain none.  The JSON certificate is common
input, but every theorem-facing entry used by this program is checked rather
than assumed.

As a further cross-check, \texttt{audit\_flint\_crosscheck.py} passes the
reconstructed integer polynomials to FLINT/Arb \cite{Johansson2017}.
Rigorous complex-ball isolation independently confirms the $41/22$ split of
the roots of $H$, one root in every exported interval, and the absence of
roots of $K$ and of the circle factors in $(-1,1)$.  This third computation is
logically redundant but guards against an implementation-specific
root-isolation error.

\subsection{Reproduction}

The supplement's README gives the complete commands and pinned dependencies.
The primary run writes a fresh JSON certificate for bytewise comparison with
the submitted one; the manuscript-consistency check and the two independent
audits are then run separately.  The archive supplies reference outputs and a
SHA-256 manifest for every file.

\section*{Data and code availability}

The manuscript source is supplied with the submission.  The exact verifier,
CAS-free independent reproduction, JSON certificate, expected outputs, file
hashes, numerical-to-exact history, exploratory transcript, and independent
FLINT/Arb root audit accompany this article as supplementary material.  The
primary verifier needs no network access and makes no floating-point proof
decision.

\section*{Funding}

This research did not receive any specific grant from funding agencies in the
public, commercial, or not-for-profit sectors.

\section*{Declaration of competing interest}

The author declares that he has no known competing financial interests or
personal relationships that could have appeared to influence the work
reported in this paper.

\section*{Declaration of generative AI and AI-assisted technologies in the
manuscript preparation process}

During the preparation of this work, the author used ChatGPT and Codex
(OpenAI) for computational exploration, generation and checking of proof
components and software, and drafting and editing the manuscript.  After
using these tools, the author reviewed and edited the resulting mathematical
arguments, computations, code, and text as needed and takes full
responsibility for the content of the article.

\begin{thebibliography}{99}

\bibitem{DraperSaad2026}
T.~L. Draper and F.~A. Saad,
\emph{Space-entropy lower bounds for random sampling},
arXiv:2607.14503 (2026).
\href{https://arxiv.org/abs/2607.14503}{arXiv:2607.14503}.

\bibitem{Zudilin2004}
W.~Zudilin,
\emph{An essay on irrationality measures of $\pi$ and other logarithms},
Chebyshevskii Sbornik \textbf{5} (2004), no.~2, 49--65.
\href{https://arxiv.org/abs/math/0404523}{arXiv:math/0404523}.

\bibitem{Rhin1987}
G.~Rhin,
\emph{Approximants de Pad\'e et mesures effectives d'irrationalit\'e},
in: C.~Goldstein (Ed.), S\'eminaire de Th\'eorie des Nombres,
Paris 1985--86, Progress in Mathematics, vol.~71,
Birkh\"auser, Boston, 1987, pp.~155--164.
\href{https://doi.org/10.1007/978-1-4757-4267-1_11}
{doi:10.1007/978-1-4757-4267-1\_11}.

\bibitem{Salikhov2007}
V.~Kh. Salikhov,
\emph{On the irrationality measure of $\log3$},
Doklady Mathematics \textbf{76} (2007), no.~3, 955--957.
\href{https://doi.org/10.1134/S1064562407060361}
{doi:10.1134/S1064562407060361}.

\bibitem{WuWang2014}
Q.~Wu and L.~Wang,
\emph{On the irrationality measure of $\log3$},
Journal of Number Theory \textbf{142} (2014), 264--273.
\href{https://doi.org/10.1016/j.jnt.2014.03.007}
{doi:10.1016/j.jnt.2014.03.007}.

\bibitem{BLS2018}
I.~V. Bondareva, M.~Yu. Luchin, and V.~Kh. Salikhov,
\emph{Symmetrized polynomials in a problem of estimating the irrationality
measure of the number $\log3$},
Chebyshevskii Sbornik \textbf{19} (2018), no.~1, 15--25.
\href{https://doi.org/10.22405/2226-8383-2018-19-1-15-25}
{doi:10.22405/2226-8383-2018-19-1-15-25}.

\bibitem{Hata1993}
M.~Hata,
\emph{Rational approximations to $\pi$ and some other numbers},
Acta Arithmetica \textbf{63} (1993), no.~4, 335--349.
\href{https://doi.org/10.4064/aa-63-4-335-349}
{doi:10.4064/aa-63-4-335-349}.

\bibitem{MontgomeryVaughan2007}
H.~L. Montgomery and R.~C. Vaughan,
\emph{Multiplicative Number Theory I: Classical Theory},
Cambridge Studies in Advanced Mathematics, vol.~97,
Cambridge University Press, Cambridge, 2007.

\bibitem{SymPy2017}
A.~Meurer, C.~P. Smith, M.~Paprocki, O.~\v{C}ert\'{\i}k,
S.~B. Kirpichev, M.~Rocklin, A.~Kumar, S.~Ivanov, J.~K. Moore,
S.~Singh, T.~Rathnayake, S.~Vig, B.~E. Granger, R.~P. Muller,
F.~Bonazzi, H.~Gupta, S.~Vats, F.~Johansson, F.~Pedregosa,
M.~J. Curry, A.~R. Terrel, \v{S}. Rou\v{c}ka, A.~Saboo,
I.~Fernando, S.~Kulal, R.~Cimrman, and A.~Scopatz,
\emph{SymPy: symbolic computing in Python},
PeerJ Computer Science \textbf{3} (2017), e103.
\href{https://doi.org/10.7717/peerj-cs.103}
{doi:10.7717/peerj-cs.103}.

\bibitem{Johansson2017}
F.~Johansson,
\emph{Arb: efficient arbitrary-precision midpoint-radius interval
arithmetic},
IEEE Transactions on Computers \textbf{66} (2017), no.~8, 1281--1292.
\href{https://doi.org/10.1109/TC.2017.2690633}
{doi:10.1109/TC.2017.2690633}.

\end{thebibliography}

\end{document}
